\documentclass{article}

% Ćirilica: kodiranje fonta i jezik (srpski, ćirilica)
\usepackage[T2A]{fontenc}
\usepackage[utf8]{inputenc}
\usepackage[serbianc]{babel}

% Veličina strane i margine
\usepackage[a4paper,top=2cm,bottom=2cm,left=3cm,right=3cm,marginparwidth=1.75cm]{geometry}

% Korisni paketi
\usepackage{amsmath}
\usepackage{graphicx}
\usepackage{float}
\usepackage{tcolorbox}
\usepackage[colorlinks=true, allcolors=blue]{hyperref}

% Цртежи (TikZ)
\usepackage{tikz}
\usetikzlibrary{arrows.meta, calc, decorations.markings, babel}

% TikZ стилови који се понављају на цртежима
\tikzset{
    naboj/.style={circle, draw, thick, minimum size=6mm, inner sep=0pt, font=\small},
    jezgro/.style={circle, draw, thick, fill=gray!15, minimum size=7mm, inner sep=0pt, font=\small},
    elektron/.style={circle, draw, fill=cyan!30, inner sep=0pt, minimum size=5pt},
    supljina/.style={circle, draw, thick, minimum size=5mm, inner sep=0pt, font=\small},
    sredstrelica/.style={postaction={decorate}, decoration={markings, mark=at position 0.55 with {\arrow{Stealth}}}},
}

\title{ОЕТ 1 - Уводно предавање}
\author{Саобраћајно техничка школа, Београд, Земун}

\begin{document}
\maketitle

%==================================================================
% БЕЛЕШКЕ ИЗ СВЕСКЕ (редом по сликама IMG_6501 -- IMG_6508)
%==================================================================

\section{Увод у разумевање струје}

% ------------------------- IMG_6501 -------------------------

\textbf{Кристијан Хајгенс}: 1629
\begin{itemize}
    \item Сат са клатном -- 1657
    \item Телескоп -- 1655 $\rightarrow$ Сатурнови прстенови
    \item Хелиоцентрични систем -- 1698 \textit{Cosmotheoros}
\end{itemize}

Кристијан Хајгенс био је Холандски математичар и физичар. Пошто је био из фамилије дипломата, отац Кристијан Хајгенс који је био композитор а пре свега секретар принчева и брат Константин Хајгенс 
који је такође био секретар принчева, и пошто је био из богате и угледне породице имао је вишак слободног времена и сигурности коју је уложио у размишљање и  откривање света око себе. Астрономијаа му је билаа страст, као и већини тадашњих мислиоца, и гледање у небо му је давало идеје којима је богатио свет око себе. 

Кроз посматрање неба открио је први Сатурнов месец, Сатурново прстење и наајважније од свега, утабаао је пут како би се развила теорија које је зачела научни метод, \textbf{Хелиоцентрични систем}.

\subsection{Атоми}

\begin{figure}[H]
    \centering
    \begin{tikzpicture}
        % језгро
        \node[jezgro] at (0,0) {Si};
        % љуске и број електрона на свакој
        \foreach \r/\n in {1/2, 1.6/8, 2.2/4} {
            \draw (0,0) circle (\r);
            \node[fill=white, inner sep=1pt, font=\small] at (\r,0) {\n};
        }
        % 1. љуска: 2 електрона
        \foreach \a in {90,270} \node[elektron] at (\a:1) {};
        % 2. љуска: 8 електрона
        \foreach \a in {22.5,67.5,...,337.5} \node[elektron] at (\a:1.6) {};
        % 3. љуска: 4 валентна електрона
        \foreach \a in {45,135,225,315} \node[elektron] at (\a:2.2) {};
        \node[font=\large] at (-2.5,2.3) {Si};
    \end{tikzpicture}
    \caption{Силицијум (Si) -- распоред електрона по љускама}
    \label{fig:si}
\end{figure}

Атоме  дефинише њихов атомски број. Атомски број представља колико је атом тежак, односно колико електрона може да држи при  себи. Атоми се распоређу по правилима дефинисаним у хемији али је најбитније да у последњем енергетском слују - \textbf{валентном} - може да се налази мааксимаално 8 електрона. Број електрона у валентном слоју дефинише проводност материјала.

По проводности, материјали се деле на: \textbf{Изолаторе(диалектрике), Проводнике и Полупроводнике(нећемо их обрађивати)}

\begin{tcolorbox}[colback=blue!5, colframe=blue!60, title=Важно]

    Више валентних електрона $\rightarrow$ изолатори.\\
    Мање валентних електрона $\rightarrow$ проводници.\\
\\
    проводници $<$ 4 валентна електрона $<$ изолатори 

\end{tcolorbox}

\begin{table}[H]
    \centering
    \begin{tabular}{|c|c|c|}
        \hline
        Елемент & Атомски број & Распоред по љускама \\
        \hline
        Si & 14 & $2+8+4$ \\
        Cu & 29 & $2+8+18+1$ \\
        Au & 79 & $2+8+18+32+18+1$ \\
        Al & 13 & $2+8+3$ \\
        Fe & 26 & $2+8+14+2$ \\
        \hline
    \end{tabular}
    \caption{Распоред електрона -- Si, Cu, Au, Al, Fe}
    \label{tab:metali}
\end{table}

\textbf{Проводници} су материјали који проводе наелектрисање са лакоћом јер атоми од којих су сачињени ти материјали немају попуњену валентну љуску.

\textbf{Изолатори} - дијалектрици - супротно, слабије проводе електроне јер атоми од којих су сачињени немају простора у својој валентној љусци да приме додатне електроне


% ------------------------- IMG_6502 -------------------------

\subsection{Историјат електрицитета}

\begin{itemize}
    \item Античко доба $\rightarrow$ трљање ћилибара, тканине
    \item 1785. Кулон -- закон електростатике
    \item 1820. А. Ампер -- открива однос електричних и магнетских појава
\end{itemize}

\subsection{Електростатика}

Тела се могу наелектрисати \textbf{без додира} -- \textbf{електростатичком индукцијом}.

\begin{equation}
    Q = \pm N e
\end{equation}

\begin{equation}
\begin{aligned}
    F &= k \frac{|Q_1| |Q_2|}{r^2},\\
    k &= 9 \cdot 10^{9} \; \frac{\mathrm{N m^2}}{\mathrm{C^2}} \quad \text{за ваздух} \\
\end{aligned}
\end{equation}

% ------------------------- IMG_6503 -------------------------

\begin{figure}[H]
    \centering
    \begin{tikzpicture}[>=Stealth]
        % 1. ред: Q1 (+), Q2 (-)
        \node[naboj, label=above:{$Q_1$}] (a1) at (0,0) {$+$};
        \node[naboj, label=above:{$Q_2$}] (b1) at (4,0) {$-$};
        \draw[->, thick] (a1.east) -- ++(1,0) node[midway, above] {$F_{21}$};
        \draw[->, thick] (b1.west) -- ++(-1,0) node[midway, above] {$F_{12}$};
        % 2. ред: Q1 (-), Q2 (-)
        \node[naboj, label=above:{$Q_1$}] (a2) at (0,-1.6) {$-$};
        \node[naboj, label=above:{$Q_2$}] (b2) at (4,-1.6) {$-$};
        \draw[->, thick] (a2.west) -- ++(-1,0) node[midway, above] {$F_{21}$};
        \draw[->, thick] (b2.east) -- ++(1,0) node[midway, above] {$F_{12}$};
        % 3. ред: Q1 (+), Q2 (+)
        \node[naboj, label=above:{$Q_1$}] (a3) at (0,-3.2) {$+$};
        \node[naboj, label=above:{$Q_2$}] (b3) at (4,-3.2) {$+$};
        \draw[->, thick] (a3.west) -- ++(-1,0) node[midway, above] {$F_{21}$};
        \draw[->, thick] (b3.east) -- ++(1,0) node[midway, above] {$F_{12}$};
    \end{tikzpicture}
    \caption{Силе између наелектрисања $Q_1$ и $Q_2$}
    \label{fig:sile}
\end{figure}

\subsection{Електрично поље}

Електрично поље је последица постојања наелектрисања. Њега дефинишемо само за ту честицу коју посматрамо.

\begin{equation}
    E = \frac{F}{Q_p}
\end{equation}

\begin{equation}
    \Rightarrow E = \frac{Q}{4 \pi \varepsilon r^2} = k\frac{Q}{r^2}
\end{equation}

\begin{figure}[H]
    \centering
    \begin{tikzpicture}[>=Stealth]
        % --- позитивно и негативно наелектрисање ---
        \node[naboj] (p) at (0,0) {$+$};
        \node[naboj] (m) at (4,0) {$-$};
        \draw[sredstrelica] (p) -- (m);
        \foreach \a in {25,50} {
            \draw[sredstrelica] (p) to[out=\a, in=180-\a] (m);
            \draw[sredstrelica] (p) to[out=-\a, in=180+\a] (m);
        }
        \foreach \a in {120,180,240} \draw[->] (p) -- ++(\a:1.3);
        \foreach \a in {60,0,-60} \draw[<-] (m) -- ++(\a:1.3);

        % --- два позитивна наелектрисања ---
        \node[naboj] (l) at (0,-4) {$+$};
        \node[naboj] (d) at (4,-4) {$+$};
        \foreach \a in {100,140,180,220,260} \draw[->] (l) -- ++(\a:1.2);
        \draw[->] (l) to[out=40, in=-100] ++(1.3,1.5);
        \draw[->] (l) to[out=-40, in=100] ++(1.3,-1.5);
        \foreach \a in {80,40,0,-40,-80} \draw[->] (d) -- ++(\a:1.2);
        \draw[->] (d) to[out=140, in=-80] ++(-1.3,1.5);
        \draw[->] (d) to[out=-140, in=80] ++(-1.3,-1.5);
    \end{tikzpicture}
    \caption{Линије електричног поља}
    \label{fig:linije-polja}
\end{figure}

% ------------------------- IMG_6504 -------------------------



\subsection{Потенцијал и напон}

\textbf{Потенцијал} -- \textbf{рад} који спољна сила мора извршити за померање пробног наелектрисања.

\begin{equation}
    V = \frac{W}{Q_p}
\end{equation}

\begin{center}
    \begin{tikzpicture}
        \fill (0,0) circle (2pt) node[above] {$V_A$};
        \fill (3,-0.3) circle (2pt) node[above] {$V_B$};
    \end{tikzpicture}
\end{center}

Са \textbf{U} дефинишемо напон. То је разлика потенцијалаа измешу 2 тачке.

\begin{equation}
    U_{AB} = V_A - V_B
\end{equation}

\subsection{Смерови електрона, струје и напона}

Посматрамо проводник (или отпорник) чији је један крај на вишем потенцијалу ($+$), а други на нижем ($-$).

\begin{itemize}
    \item \textbf{Кретање електрона одражава се као скакање електронаа са једног атома на други.}
    Кретање електрона за собом оставља празан енергетски простор -- \textbf{шупљине}, које се привидно крећу у супротном смеру.
\end{itemize}

\begin{figure}[H]
    \centering
    \begin{tikzpicture}[>=Stealth]
        \draw[->, thick] (-0.4,0) -- (-1.9,0);
        \node[supljina] at (0,0) {};
        \node[supljina] at (1.2,0) {$e$};
        \draw[->, thick] (1.6,0) -- (3.1,0);
    \end{tikzpicture}
    \caption{Шупљина (лево) и електрон (десно) се крећу у супротним смеровима}
    \label{fig:supljina}
\end{figure}

\begin{itemize}
    \item Струја представља усмерено кретање \textbf{шупљина}.
\end{itemize}

\begin{itemize}
    \item Кретање електрона проузрокује електрично поље у супротном смеру.
\end{itemize}

\begin{figure}[H]
    \centering
    \begin{tikzpicture}[>=Stealth]
        \node at (0,0.5) {$e$};
        \draw[->, thick] (0.3,0.5) -- (2.8,0.5);
        \node at (0,0) {$U$};
        \draw[->, thick] (2.8,0) -- (0.3,0);
    \end{tikzpicture}
    \caption{Смер кретања електрона ($e$) и напона ($U$)}
    \label{fig:e-u}
\end{figure}

Да се подсетимо. Електрон се креће у једном смеру па се струја и напон крећу у супротном.

\begin{figure}[H]
    \centering
    \begin{tikzpicture}[>=Stealth]
        \node at (-0.6,0) {$-$};
        \node at (5.6,0) {$+$};
        % електрон: креће се ка +
        \node[circle, draw, thick, minimum size=6mm, inner sep=0pt] (e) at (0.3,0.9) {};
        \node[above left] at (e.north) {$e^-$};
        \draw[->, thick] (e.east) -- (4.6,0.9);
        % шупљина / струја: креће се ка -
        \node[circle, draw, thick, minimum size=6mm, inner sep=0pt] (h) at (4.7,0) {};
        \draw[->, thick] (h.west) -- (0.4,0) node[midway, below] {$I$};
        % смер напона
        \draw[->, thick] (4.9,-1) -- (0.4,-1) node[midway, below] {$V$};
    \end{tikzpicture}
    \caption{Електрон се креће од $-$ ка $+$, а струја $I$ и смер напона $V$ од $+$ ка $-$}
    \label{fig:smerovi}
\end{figure}

АЛИ, ОВО ЈЕ ЈАКО ВАЖНО.

\begin{tcolorbox}[colback=blue!5, colframe=blue!60, title=ЈАКО ВАЖНО]

Смер кретања електричног поља се на шемама означава стрелицом која је оријентисаана \textbf{СУПРОТНО} од смера кретања напона са ознаком + (плус) на краају те стрелице.

\end{tcolorbox}

\begin{figure}[H]
    \centering
    \begin{tikzpicture}[>=Stealth]
        % грана кола
        % смер напона: од + ка -
        \draw[->, thick] (4.4,0.5) -- (0.6,0.5) node[midway, above] {$V$};
        % ознака напона: показује на +
        \draw[->, thick] (0.2,1) to[out=70, in=110] (4.8,0.75);
        \node at (2.5,2.35) {$V$};
        \node at (5.25,0.85) {$+$};
    \end{tikzpicture}
    \caption{Права стрелица: смер напона (од $+$ ка $-$). Закривљена стрелица: ознака напона, показује на $+$}
    \label{fig:oznaka-napona}
\end{figure}

Овако изгледа на примеру отпорника, пошто ћемо овај случај виђати јако често.

\begin{figure}[H]
    \centering
    \begin{tikzpicture}[>=Stealth]
        % отпорник између два краја
        \draw[thick] (0,0) circle (2pt) -- (1.2,0)
            -- (1.35,0.35) -- (1.65,-0.35) -- (1.95,0.35) -- (2.25,-0.35) -- (2.55,0.35) -- (2.7,0)
            -- (3.9,0) circle (2pt);
        % ознака напона: показује на +
        \draw[->, thick] (2.9,0.8) to[out=110, in=60] (0.15,0.65);
        \node at (1.6,1.85) {$V$};
        \node at (-0.3,0.45) {$+$};
        % струја: од + ка -
        \draw[->, thick] (0.4,-0.8) -- (3.5,-0.8) node[midway, below] {$I$};
    \end{tikzpicture}
    \caption{Отпорник: струја $I$ тече од $+$ ка $-$, а ознака напона $V$ показује на $+$}
    \label{fig:otpornik-napon}
\end{figure}

\begin{tcolorbox}[colback=blue!5, colframe=blue!60, title=ВАЖНО ЗА ЗАПАМТИТИ]
\begin{itemize}
    \item \textbf{Електрони} ($e^-$) се крећу од $-$ ка $+$, јер их привлачи позитиван крај.
    \item \textbf{Струја} $I$ (технички смер) тече од $+$ ка $-$, тј. \textbf{супротно} од електрона. То је исти смер у коме се привидно крећу шупљине.
    \item \textbf{Смер напона} $V$ (смер пада потенцијала) је од $+$ ка $-$, исто као струја кроз отпорник.
    \item \textbf{Ознака напона} (стрелица коју цртамо у колу) се црта \textbf{супротно} од смера напона: креће од $-$ и \textbf{показује на $+$}. Она је дакле у истом смеру у коме се крећу електрони.
\end{itemize}
\end{tcolorbox}

\subsection{Електростатичка индукција}

Електростатичка индукција представља наелектрисавање тела или честица, без додира или трљања.

До сада смо посматрали како се понаша једна наелектрисана честица. 
Уколико не посматрамо више само тачкаста наелектрисања, услед електростатичке индукције која делује на неко тело, наелектрисање се гомила по површини најближој спољном пољу.

\begin{figure}[H]
    \centering
    \begin{tikzpicture}[>=Stealth]
        \node[naboj] (q) at (0,0) {$+$};
        \draw[->] (q) to[out=60, in=180] (2.2,1.1);
        \draw[->] (q) to[out=-30, in=180] (2.2,-0.6);
        % површина тела
        \draw[thick] (2.5,-1.8) -- (2.5,1.8);
        \foreach \y in {1.4,1.0,...,-1.4} \node at (2.8,\y) {$-$};
    \end{tikzpicture}
    \caption{Електростатичка индукција}
    \label{fig:indukcija}
\end{figure}

% TODO: нисам успео да прочитам заокружену реч испод „Пример:“ -- провери у свесци.
Пример: \textbf{громобран}, \textbf{Рашије}

\begin{figure}[H]
    \centering
    \begin{tikzpicture}[>=Stealth]
        % громобран
        \draw[thick] (-0.9,0) -- (0,3) -- (0.9,0);
        \node[naboj, minimum size=5mm] at (0,1.75) {$-$};
        \node[naboj, minimum size=5mm] at (0.6,3.6) {$+$};
        \node at (0,4.4) {громобран};
        % облак и тло
        \begin{scope}[xshift=4cm]
            \draw[thick] (-1.3,3) -- (1.3,3);
            \foreach \x in {-0.8,0,0.8} \node at (\x,3.3) {$+$};
            \draw[->] (-0.6,1.6) -- (-0.6,2.7);
            \draw[->] (0.6,2.7) -- (0.6,1.6);
            \draw[thick] (-0.9,0) -- (0,1.3) -- (0.9,0);
            \node at (0,0.5) {$-$};
        \end{scope}
    \end{tikzpicture}
    \caption{Пример: громобран}
    \label{fig:gromobran}
\end{figure}

% ------------------------- IMG_6508 -------------------------

\subsection{Фарадејев кавез}

\begin{figure}[H]
    \centering
    \begin{tikzpicture}[>=Stealth]
        \node[naboj] (q) at (0,0) {$+$};
        % зидови кавеза
        \draw[thick] (1.6,-1.5) rectangle (5.1,1.5);
        \draw[thick] (2.0,-1.1) rectangle (4.7,1.1);
        % индуковано наелектрисање
        \foreach \y in {0.85,0.3,-0.25,-0.8} \node[naboj, minimum size=3.5mm, font=\scriptsize] at (1.8,\y) {$-$};
        \foreach \x in {2.3,3.0,3.7,4.4} \node[naboj, minimum size=3.5mm, font=\scriptsize] at (\x,1.3) {$-$};
        \node at (3.35,0) {$E = 0$};
        \draw[->] (q) to[out=50, in=180] (1.55,0.9);
        \draw[->] (q) to[out=-50, in=180] (1.55,-0.9);
    \end{tikzpicture}
    \caption{Фарадејев кавез}
    \label{fig:faradej}
\end{figure}

\section{ДОДАТАК ЗА ЗАИНТЕРЕСОВАНЕ}

\subsection{Електрично поље у диелектрику}

Када се изолатор(диалектрик) нађе у електричном пољу, његови молекули се \textbf{поларизују}: постају диполи окренути у смеру поља. Унутар изолатора се поље поништава, па некомпензована наелектрисања остају само на површини изолатора. Та \textbf{поларизациона наелектрисања} стварају поље $E_p$ које је супротног смера од поља $E_0$ које би наелектрисане плоче стварале у вакууму (ваздуху). Зато је поље у изолатору \textbf{слабије}:

\begin{equation}
    E = E_0 - E_p
\end{equation}

Смањење поља може се приказати и преко релативне диелектричне константе $\varepsilon_r$:

\begin{equation}
    E = \frac{E_0}{\varepsilon_r}
\end{equation}

За вакуум је $\varepsilon_r = 1$, а за све остале изолаторе $\varepsilon_r > 1$. Што се материјал лакше поларизује, то је $\varepsilon_r$ веће, а поље у њему слабије. На пример, у води је поље 80 пута слабије него у вакууму. \textbf{Апсолутна диелектрична константа} је:

\begin{equation}
    \varepsilon = \varepsilon_r \cdot \varepsilon_0, \qquad \varepsilon_0 = 8{,}85 \cdot 10^{-12} \; \frac{\mathrm{F}}{\mathrm{m}}
\end{equation}

Јединица $\mathrm{F/m}$ (фарад по метру) је исто што и $\mathrm{C^2/(N\,m^2)}$.

\begin{table}[H]
    \centering
    \begin{tabular}{|l|c|c|}
        \hline
        \textbf{Диелектрик} & $\varepsilon_r$ & $\varepsilon = \varepsilon_r \varepsilon_0$ \; [$10^{-12}$ F/m] \\
        \hline
        вакуум & 1 & 8,85 \\
        ваздух & 1,00059 & $\approx$ 8,85 \\
        стакло & 2--16 & 17,7--141,6 \\
        лискун & 4--8 & 35,4--70,8 \\
        порцелан & 4,5--6 & 39,8--53,1 \\
        вода & 80 & 708 \\
        \hline
        \multicolumn{3}{|l|}{\textit{\textbf{Остали често коришћени изолатори}}} \\
        \hline
        гума (природна) & 3--7 & 26,5--62,0 \\
        силиконска гума & 2,9--4 & 25,7--35,4 \\
        неопрен & 6,7 & 59,3 \\
        PVC (изолација каблова) & 3,6--6 & 31,9--53,1 \\
        полиетилен & 2,25 & 19,9 \\
        тефлон (PTFE) & 2,1 & 18,6 \\
        полистирен & 2,4--2,7 & 21,2--23,9 \\
        папир & 2,3 & 20,4 \\
        трансформаторско уље & 2,2 & 19,5 \\
        бакелит & 3,5--5 & 31,0--44,3 \\
        суво дрво & 2--6 & 17,7--53,1 \\
        најлон & 4--5 & 35,4--44,3 \\
        FR-4 (штампане плочице) & 4,5 & 39,8 \\
        $\mathrm{SiO_2}$ (оксид у чиповима) & 3,9 & 34,5 \\
        \hline
        \multicolumn{3}{|l|}{\textit{\textbf{Полупроводник (за PN спојеве})}} \\
        \hline
        силицијум (Si) & 11,7 & 103,5 \\
        \hline
    \end{tabular}
    \caption{Релативне и апсолутне диелектричне константе неких материјала}
    \label{tab:dielektrici}
\end{table}

\subsection{Водоник, угљеник, кисеоник}

Проводност ових атома можемо и сликовито приказати. Видимо да се у случају кисеоника, у ваалентној љуспи(зони) наалаази по 6 електронаа, што га чини слабо проводним. Односно ближе изолаатору него проводнику.

\begin{table}[H]
    \centering
    \begin{tabular}{|c|c|c|}
        \hline
        Елемент & Атомски број & Распоред по љускама \\
        \hline
        H & 1 & $1$ \\
        C & 6 & $2+4$ \\
        O & 8 & $2+6$ \\
        \hline
    \end{tabular}
    \caption{Распоред електрона -- H, C, O}
    \label{tab:hco}
\end{table}

У случају воде, могли бисмо да кажемо исту ствар. Додуше овде се јавља једна цака. Вода је молекул, не само атом. Такође је у течном стању. По законима хемије  у које нећу залазити, чини је бољим проводником од кисеоника, али и даље недовољно добрим.

\begin{figure}[H]
    \centering
    \begin{tikzpicture}
        \coordinate (L) at (0,0);
        \coordinate (R) at (2.6,0);
        \foreach \c in {L,R} {
            \node[jezgro] at (\c) {O};
            \draw (\c) circle (0.6);   % 1. љуска
            \draw (\c) circle (1.5);   % 2. љуска
        }
        % 1. љуска -- по 2 електрона
        \node[elektron] at ($(L)+(90:0.6)$) {};
        \node[elektron] at ($(L)+(250:0.6)$) {};
        \node[elektron] at ($(R)+(70:0.6)$) {};
        \node[elektron] at ($(R)+(270:0.6)$) {};
        % 2. љуска -- електрони који нису у вези
        \foreach \a in {105,130,160,220} \node[elektron] at ($(L)+(\a:1.5)$) {};
        \foreach \a in {5,-15,-45,-88} \node[elektron] at ($(R)+(\a:1.5)$) {};
        % заједнички електрони (у пресеку две љуске)
        \node[elektron] at (1.3,0.75) {};
        \node[elektron] at (1.3,-0.75) {};
        \node[font=\large] at (1.3,2.1) {$\mathrm{O_2}$};
    \end{tikzpicture}
    \caption{Кисеоник $\mathrm{O_2}$}
    \label{fig:o2}
\end{figure}

\begin{figure}[H]
    \centering
    \begin{tikzpicture}
        \coordinate (O) at (0,0);
        \coordinate (HL) at (-2.3,0);
        \coordinate (HR) at (2.3,0.1);
        % кисеоник
        \node[jezgro] at (O) {O};
        \draw (O) circle (0.6);
        \draw (O) circle (1.5);
        \node[elektron] at ($(O)+(115:0.6)$) {};
        \node[elektron] at ($(O)+(-40:0.6)$) {};
        \foreach \a in {120,78,-82,-39} \node[elektron] at ($(O)+(\a:1.5)$) {};
        % водоници
        \foreach \c in {HL,HR} {
            \node[jezgro] at (\c) {H};
            \draw (\c) circle (0.85);
        }
        % заједнички електрони (H -- O)
        \node[elektron] at (-1.47,0) {};
        \node[elektron] at (1.47,0.06) {};
        \node[font=\large] at (0,2.1) {$\mathrm{H_2O}$};
    \end{tikzpicture}
    \caption{Вода $\mathrm{H_2O}$}
    \label{fig:h2o}
\end{figure}

\subsection{Молекул гуме}
\textbf{Молекул изопрена} $\mathrm{C_5H_8}$ (основна јединица гуме):

\begin{figure}[H]
    \centering
    \begin{tikzpicture}[scale=0.75,
        jez/.style={circle, draw, thick, fill=gray!15, minimum size=6mm, inner sep=0pt, font=\small},
        el/.style={circle, draw, fill=cyan!30, inner sep=0pt, minimum size=5pt}]
        % --- угљеници: 1. љуска (2 e) и 2. љуска (валентна) ---
        \foreach \x/\y in {0/0, 2.5/0, 5/0, 7.5/0, 2.5/2.5} {
            \draw (\x,\y) circle (0.55);
            \draw (\x,\y) circle (1.5);
            \node[el] at ($(\x,\y)+(45:0.55)$) {};
            \node[el] at ($(\x,\y)+(225:0.55)$) {};
            \node[jez] at (\x,\y) {C};
        }
        % --- водоници: само једна љуска ---
        \foreach \x/\y in {-1.41/1.41, -1.41/-1.41, 0.5/2.5, 2.5/4.5, 4.5/2.5, 5/-2, 8.91/1.41, 8.91/-1.41} {
            \draw (\x,\y) circle (0.85);
            \node[jez] at (\x,\y) {H};
        }
        % --- заједнички електрони (у пресеку две љуске) ---
        % двоструке везе C=C: по 4 електрона
        \foreach \x/\y in {1.25/0.2, 1.25/-0.2, 1.25/0.55, 1.25/-0.55,
                           6.25/0.2, 6.25/-0.2, 6.25/0.55, 6.25/-0.55}
            \node[el] at (\x,\y) {};
        % једноструке везе C-C: по 2 електрона
        \foreach \x/\y in {3.75/0.2, 3.75/-0.2, 2.3/1.25, 2.7/1.25}
            \node[el] at (\x,\y) {};
        % везе C-H: по 2 електрона
        \foreach \x/\y in {-1.08/0.8, -0.8/1.08, -0.8/-1.08, -1.08/-0.8,
                           1.18/2.3, 1.18/2.7, 2.3/3.83, 2.7/3.83, 3.83/2.7, 3.83/2.3,
                           5.2/-1.32, 4.8/-1.32, 8.3/1.08, 8.58/0.8, 8.58/-0.8, 8.3/-1.08}
            \node[el] at (\x,\y) {};
    \end{tikzpicture}
    \caption{Изопрен $\mathrm{C_5H_8}$ -- заједнички електрони у пресецима љуски}
    \label{fig:izopren}
\end{figure}

Сваки угљеник има 2 електрона у првој љусци и 8 у другој, а сваки водоник 2. Сви валентни електрони су у везама. У пресеку две љуске угљеника \textbf{2 електрона} значе једноструку везу, а \textbf{4 електрона} двоструку везу. Пети угљеник није у низу са осталима, већ стоји изнад другог као бочна група $\mathrm{CH_3}$.

Пошто су сви валентни електрони попуњени, ово чини гуму идеаалним изолатором. Не може ни један електрон да уђе или да прође кроз молекуле гуме :).

\end{document}